problem:  
Let \( (X,J,\omega,\Omega,H) \) be a compact complex threefold satisfying the Hull–Strominger system, and let \(\widehat H \in \widehat H^3(X,\mathbb{Z})\) denote the quantized flux, whose cohomology component \([H]\in H^3(X,\mathbb{Z})\) may include torsion.  
De la Ossa–Svanes constructed a differential graded Lie algebra \(\mathcal{L}_{X,H}\) controlling analytic deformations of the Strominger system. In the formal moduli framework of Lurie–Pridham, each such dg Lie algebra defines a formal moduli problem \(\mathfrak{Def}(\mathcal{L}_{X,H})\).  

Let \(\mathcal{L}_{X,H}^{\mathrm{fix}}\) denote the deformation dg Lie algebra obtained by **restricting infinitesimal variations** \((\delta J,\delta\omega,\delta F,\delta H)\) so that the associated differential cohomology class \(\widehat H\) remains fixed in \(\widehat H^3(X,\mathbb{Z})\).  This constraint can be modeled by taking the **homotopy fiber** of the natural morphism  
\[
\mathcal{L}_{X,H}\longrightarrow \mathrm{Tot}\,\Omega^{3}_{\mathrm{cl}}(X),
\]
composed with the projection to the integral lattice in differential cohomology.

> **Problem:**  
> Given two fluxes \(\widehat H_1, \widehat H_2 \in \widehat H^3(X,\mathbb{Z})\) with identical de Rham component but distinct torsion classes in \(H^3(X,\mathbb{Z})_{\mathrm{tors}}\):  
> 1. Determine whether the dg Lie algebras \(\mathcal{L}_{X,H_1}^{\mathrm{fix}}\) and \(\mathcal{L}_{X,H_2}^{\mathrm{fix}}\) are quasi‑isomorphic.  
> 2. Equivalently, does imposing the differential‑cohomological flux constraint cause the derived formal moduli spaces \(\mathfrak{Def}(\mathcal{L}_{X,H_i}^{\mathrm{fix}})\) to depend on the torsion component of \([H_i]\)?  
> 3. Can torsion therefore introduce **formally inequivalent deformation sectors** of the Strominger system that are invisible to the de Rham deformation complex?  

Why is it a "good" problem:  
This problem is now sharply formulated: it works entirely within the **formal moduli problem** framework given by the Lurie–Pridham correspondence, using the established de la Ossa–Svanes dg‑Lie algebra, and introduces the **homotopy‑fiber construction** to encode the integral flux constraint rigorously. It asks whether torsion in differential cohomology modifies the quasi‑isomorphism type of the deformation dg‑Lie algebra—precisely formulated and mathematically meaningful.  
The question probes how **discrete topological data** (torsion in \([H]\)) could alter the **derived deformation theory** of heterotic compactifications—an intersection of **non‑Kähler geometry**, **derived deformation theory**, and **differential topology**. If torsion can be detected at the dg‑Lie or formal‑moduli level, it would reveal a new coupling between integral structure and analytic deformation theory, opening a bridge between **differential cohomology** and **derived geometric analysis**.  
It is conceptually simple, rigorously posed, logically sound, and the answer is entirely unknown—making it a genuinely “good’’ mathematical research problem.